\documentclass{amsart}
% For dealing with references we use the comment environment
\usepackage{verbatim}
\newenvironment{reference}{\comment}{\endcomment}
%\newenvironment{reference}{}{}
\newenvironment{slogan}{\comment}{\endcomment}
\newenvironment{history}{\comment}{\endcomment}

% For commutative diagrams we use Xy-pic
\usepackage[all]{xy}

% We use 2cell for 2-commutative diagrams.
\xyoption{2cell}
\UseAllTwocells

% We use multicol for the list of chapters between chapters
\usepackage{multicol}

% This is generally recommended for better output
\usepackage{lmodern}
\usepackage[T1]{fontenc}

% For cross-file-references

% Package for hypertext links:
\usepackage{hyperref}

% For any local file, say "hello.tex" you want to link to please

% Theorem environments.
%
\theoremstyle{plain}
\newtheorem{theorem}[subsection]{Theorem}
\newtheorem{proposition}[subsection]{Proposition}
\newtheorem{lemma}[subsection]{Lemma}

\theoremstyle{definition}
\newtheorem{definition}[subsection]{Definition}
\newtheorem{example}[subsection]{Example}
\newtheorem{exercise}[subsection]{Exercise}
\newtheorem{situation}[subsection]{Situation}

\theoremstyle{remark}
\newtheorem{remark}[subsection]{Remark}
\newtheorem{remarks}[subsection]{Remarks}

\numberwithin{equation}{subsection}

% Macros
%
\def\lim{\mathop{\mathrm{lim}}\nolimits}
\def\colim{\mathop{\mathrm{colim}}\nolimits}
\def\Spec{\mathop{\mathrm{Spec}}}
\def\Hom{\mathop{\mathrm{Hom}}\nolimits}
\def\Ext{\mathop{\mathrm{Ext}}\nolimits}
\def\SheafHom{\mathop{\mathcal{H}\!\mathit{om}}\nolimits}
\def\SheafExt{\mathop{\mathcal{E}\!\mathit{xt}}\nolimits}
\def\Sch{\mathit{Sch}}
\def\Mor{\mathop{\mathrm{Mor}}\nolimits}
\def\Ob{\mathop{\mathrm{Ob}}\nolimits}
\def\Sh{\mathop{\mathit{Sh}}\nolimits}
\def\NL{\mathop{N\!L}\nolimits}
\def\CH{\mathop{\mathrm{CH}}\nolimits}
\def\proetale{{pro\text{-}\acute{e}tale}}
\def\etale{{\acute{e}tale}}
\def\QCoh{\mathit{QCoh}}
\def\Ker{\mathop{\mathrm{Ker}}}
\def\Im{\mathop{\mathrm{Im}}}
\def\Coker{\mathop{\mathrm{Coker}}}
\def\Coim{\mathop{\mathrm{Coim}}}

% Boxtimes
%
\DeclareMathSymbol{\boxtimes}{\mathbin}{AMSa}{"02}

%
% Macros for moduli stacks/spaces
%
\def\QCohstack{\mathcal{QC}\!\mathit{oh}}
\def\Cohstack{\mathcal{C}\!\mathit{oh}}
\def\Spacesstack{\mathcal{S}\!\mathit{paces}}
\def\Quotfunctor{\mathrm{Quot}}
\def\Hilbfunctor{\mathrm{Hilb}}
\def\Curvesstack{\mathcal{C}\!\mathit{urves}}
\def\Polarizedstack{\mathcal{P}\!\mathit{olarized}}
\def\Complexesstack{\mathcal{C}\!\mathit{omplexes}}
% \Pic is the operator that assigns to X its picard group, usage \Pic(X)
% \Picardstack_{X/B} denotes the Picard stack of X over B
% \Picardfunctor_{X/B} denotes the Picard functor of X over B
\def\Pic{\mathop{\mathrm{Pic}}\nolimits}
\def\Picardstack{\mathcal{P}\!\mathit{ic}}
\def\Picardfunctor{\mathrm{Pic}}
\def\Deformationcategory{\mathcal{D}\!\mathit{ef}}

\setlength{\emergencystretch}{3em}
\begin{document}
\title[Stacks Project, Tag 0DHV]{404 --- Relative perfectness and finite complexes of flat modules}
\maketitle
\noindent Copyright (C) 2005--2025 Johan de Jong. Stacks Project authors. Source: \href{https://stacks.math.columbia.edu/tag/0DHV}{Tag 0DHV}.

\noindent Editorial difficulty note (not author text). Difficulty H2: The reverse finiteness implication descends the algebra and flat module to a Noetherian finite-type model, resolves there, and checks that base change preserves exactness. The other direction truncates a finite-free representative at a Tor-amplitude bound, identifying a flat final cokernel..

\noindent Editorial provenance note (not author text). History: excerpt from revision \texttt{a0444\allowbreak 6e57e\allowbreak c1fbc\allowbreak 252a8\allowbreak 71afc\allowbreak ec775\allowbreak 2fb28\allowbreak 07b14}, more-algebra.tex, lines 21979--22019. Reference presentation was adapted; original-author context and core support blocks were retained or added. The mathematical wording is unchanged. Licensed under GNU FDL 1.2 or later; no invariant sections or cover texts. See ../licenses/COPYING. Context blocks reproduce author definitions and supporting results; they are not additional counted problems.

\begin{definition}
\label{definition-relatively-perfect}
Let $R \to A$ be a flat ring map of finite presentation.
An object $K$ of $D(A)$ is {\it $R$-perfect} or {\it perfect relative to $R$}
if $K$ is pseudo-coherent
(Definition \ref{definition-pseudo-coherent})
and has finite tor dimension over $R$
(Definition \ref{definition-tor-amplitude}).
\end{definition}

\begin{definition}
\label{definition-pseudo-coherent}
Let $R$ be a ring. Denote $D(R)$ its derived category.
Let $m \in \mathbf{Z}$.
\begin{enumerate}
\item An object $K^\bullet$ of $D(R)$ is {\it $m$-pseudo-coherent}
if there exists a bounded complex $E^\bullet$ of finite free $R$-modules
and a morphism $\alpha : E^\bullet \to K^\bullet$ such that
$H^i(\alpha)$ is an isomorphism for $i > m$ and $H^m(\alpha)$
is surjective.
\item An object $K^\bullet$ of $D(R)$ is {\it pseudo-coherent}
if it is quasi-isomorphic to a bounded above complex of finite
free $R$-modules.
\item An $R$-module $M$ is called {\it $m$-pseudo-coherent}
if $M[0]$ is an $m$-pseudo-coherent object of $D(R)$.
\item An $R$-module $M$ is called
{\it pseudo-coherent}\footnote{This clashes with what is meant by
a pseudo-coherent module in \href{https://stacks.math.columbia.edu/bibliography}{Bibliography: Bourbaki-CA}.}
if $M[0]$ is a pseudo-coherent object of $D(R)$.
\end{enumerate}
\end{definition}

\begin{definition}
\label{definition-tor-amplitude}
Let $R$ be a ring. Denote $D(R)$ its derived category.
Let $a, b \in \mathbf{Z}$.
\begin{enumerate}
\item An object $K^\bullet$ of $D(R)$ has
{\it tor-amplitude in $[a, b]$}
if $H^i(K^\bullet \otimes_R^\mathbf{L} M) = 0$ for all $R$-modules
$M$ and all $i \not \in [a, b]$.
\item An object $K^\bullet$ of $D(R)$ has {\it finite tor dimension}
if it has tor-amplitude in $[a, b]$ for some $a, b$.
\item An $R$-module $M$ has {\it tor dimension $\leq d$}
if $M[0]$ as an object of $D(R)$ has tor-amplitude in $[-d, 0]$.
\item An $R$-module $M$ has {\it finite tor dimension}
if $M[0]$ as an object of $D(R)$ has finite tor dimension.
\end{enumerate}
\end{definition}

\begin{lemma}
\label{lemma-structure-relatively-perfect}
Let $R \to A$ be a flat ring map of finite presentation.
Let $K \in D(A)$. The following are equivalent
\begin{enumerate}
\item $K$ is $R$-perfect, and
\item $K$ is isomorphic to a finite complex of $R$-flat,
finitely presented $A$-modules.
\end{enumerate}
\end{lemma}

\begin{proof}
To prove (2) implies (1) it suffices by
Lemma \href{https://stacks.math.columbia.edu/tag/0DHT}{lemma cone relatively perfect (Tag 0DHT)}
to show that an $R$-flat, finitely presented $A$-module $M$ defines
an $R$-perfect object of $D(A)$. Since $M$ has finite tor dimension
over $R$, it suffices to show that $M$ is pseudo-coherent. By
Algebra, Lemma \href{https://stacks.math.columbia.edu/tag/02JO}{lemma flat finite presentation limit flat (Tag 02JO)}
there exists a finite type $\mathbf{Z}$-algebra $R_0 \subset R$
and a flat finite type ring map $R_0 \to A_0$ and
a finite $A_0$-module $M_0$ flat over $R_0$ such that
$A = A_0 \otimes_{R_0} R$ and $M = M_0 \otimes_{R_0} R$. By
Lemma \href{https://stacks.math.columbia.edu/tag/066E}{lemma Noetherian pseudo coherent (Tag 066E)}
we see that $M_0$ is pseudo-coherent $A_0$-module.
Choose a resolution $P_0^\bullet \to M_0$ by finite free
$A_0$-modules $P_0^n$. Since $A_0$ is flat over $R_0$,
this is a flat resolution. Since $M_0$ is flat over $R_0$
we find that $P^\bullet = P_0^\bullet \otimes_{R_0} R$
still resolves $M = M_0 \otimes_{R_0} R$. (You can use
Lemma \href{https://stacks.math.columbia.edu/tag/0661}{lemma base change comparison (Tag 0661)} to see this.)
Hence $P^\bullet$ is a finite free resolution of $M$
over $A$ and we conclude that $M$ is pseudo-coherent.

\medskip\noindent
Assume (1). We can represent $K$ by a bounded above complex
$P^\bullet$ of finite free $A$-modules. Assume that
$K$ viewed as an object of $D(R)$ has tor amplitude in $[a, b]$.
By Lemma \href{https://stacks.math.columbia.edu/tag/0653}{lemma last one flat (Tag 0653)} we see that
$\tau_{\geq a}P^\bullet$ is a complex of $R$-flat, finitely
presented $A$-modules representing $K$.
\end{proof}
\end{document}
